\chapter{Étude Préalable}
\label{chap:etude-prealable}

\section*{Introduction du chapitre}
\addcontentsline{toc}{section}{Introduction du chapitre}

Ce premier chapitre a pour objectif de poser le cadre général du projet. Il permet de comprendre le contexte dans lequel s'inscrit la solution proposée, d'analyser les limites des pratiques actuelles, et d'identifier précisément les besoins auxquels le système devra répondre.

Nous y présentons également la démarche méthodologique adoptée pour la réalisation du projet, basée sur une approche agile, mieux adaptée à la nature évolutive et expérimentale des solutions intégrant de l'intelligence artificielle.


\section{Contexte général du projet}
\label{sec:contexte-general}

\subsection{Cadre du projet}

Ce projet s'inscrit dans le cadre d'un projet de fin d'études, réalisé au sein de la société Tritux. Il a pour objectif principal la conception et le développement d'une plateforme intelligente de gestion des ressources humaines.

L'ambition n'est pas uniquement de digitaliser les processus existants, mais de les repenser en profondeur en intégrant des mécanismes d'automatisation avancée et d'intelligence artificielle, afin d'améliorer l'efficacité globale des équipes RH.

La durée totale du projet s'étend sur plusieurs mois, avec une organisation en sprints bihebdomadaires permettant une validation progressive des fonctionnalités et une adaptation rapide aux besoins.


\subsection{Présentation de l'organisme d'accueil}

Tritux est une entreprise spécialisée dans le domaine des technologies de l'information, proposant des solutions innovantes adaptées aux besoins des entreprises. L'entreprise intervient sur plusieurs domaines :

\begin{itemize}
    \item Le développement d'applications web et mobiles,
    \item L'intégration de systèmes \acrfull{api} REST et microservices,
    \item L'automatisation des workflows via des outils d'orchestration,
    \item L'exploitation de services d'intelligence artificielle pour l'analyse et la structuration de données,
    \item L'infrastructure cloud et la containerisation Docker.
\end{itemize}

Dans un contexte de croissance et de structuration interne, l'entreprise fait face à des problématiques liées à la gestion des profils internes et au traitement des candidatures externes, ce qui a motivé la mise en place de ce projet.


\subsection{Problématique}
\label{subsec:problematique}

Malgré la disponibilité de nombreux outils RH, une grande partie des processus reste encore manuelle, fragmentée et peu optimisée. Parmi les principales difficultés rencontrées :

\begin{itemize}
    \item \textbf{L'absence de standardisation des CV}, rendant leur lecture et leur comparaison complexes;
    
    \item \textbf{La perte de temps} liée à la saisie et à la modification manuelle des informations;
    
    \item \textbf{Le manque de centralisation} des données entre les profils internes et les candidats externes;
    
    \item \textbf{L'absence d'outils intelligents} capables d'assister les RH dans l'optimisation des profils;
    
    \item \textbf{Une gestion limitée} des candidatures et des communications associées.
\end{itemize}

Ces limitations montrent clairement que l'automatisation seule ne suffit plus. Il devient nécessaire de concevoir des systèmes capables non seulement d'exécuter des tâches, mais aussi de comprendre, structurer et valoriser l'information.


\section{Description de l'existant}
\label{sec:description-existant}

\subsection{Étude de l'existant}
\label{subsec:etude-existant}

Actuellement, la gestion des ressources humaines repose sur une combinaison d'outils classiques tels que les emails, les fichiers bureautiques (CV en formats variés) et parfois des plateformes de recrutement externes.

Les CV sont généralement traités manuellement, modifiés individuellement, et stockés sans structure uniforme. De plus, la gestion des candidatures externes est souvent dissociée de la gestion des profils internes. Les équipes RH doivent jongler entre plusieurs outils et sources de données, ce qui complique le suivi et l'analyse des candidats.

Le processus global se décompose comme suit :

\begin{enumerate}
    \item Réception des candidatures par email ou formulaires web,
    \item Tri manuel et classification,
    \item Extraction manuelle des informations pertinentes,
    \item Adaptation du CV aux standards internes,
    \item Évaluation et classement des candidats,
    \item Communication avec les candidats (emails, appels).
\end{enumerate}


\subsection{Critique de l'existant}
\label{subsec:critique-existant}

Cette approche présente plusieurs limites significatives :

\begin{itemize}
    \item \textbf{Processus long et répétitif} : chaque candidature requiert une intervention manuelle à plusieurs étapes;
    
    \item \textbf{Forte dépendance aux interventions humaines} : le succès du traitement dépend largement de l'expérience et de l'attention de chacun;
    
    \item \textbf{Risque élevé d'erreurs ou d'incohérences} : les données peuvent être mal interprétées ou mal saisies;
    
    \item \textbf{Difficulté de mise à l'échelle} : le processus devient rapidement saturé en cas de volume important de candidatures;
    
    \item \textbf{Absence d'intelligence} dans le traitement des données : aucune aide pour prioriser ou optimiser les profils.
\end{itemize}

En particulier, l'absence de standardisation des CV constitue un frein majeur à leur exploitation efficace. Chaque candidat utilise un format différent, ce qui rend difficile l'extraction cohérente des informations.


\subsection{Solution proposée}
\label{subsec:solution-proposee}

Pour répondre à ces limitations, la solution proposée consiste en une plateforme centralisée et intelligente permettant :

\begin{itemize}
    \item \textbf{L'importation et l'analyse automatique} des CV via des modèles d'intelligence artificielle;
    
    \item \textbf{L'extraction des informations} à partir des documents (formation, expérience, compétences, etc.);
    
    \item \textbf{La modification et l'optimisation} des profils, manuellement ou via un assistant intelligent;
    
    \item \textbf{La transformation des CV} selon des templates personnalisables respectant la charte de l'entreprise;
    
    \item \textbf{La gestion des offres d'emploi} synchronisées avec un site web public;
    
    \item \textbf{Le classement automatique} des candidatures via un système de scoring intelligent;
    
    \item \textbf{L'automatisation des communications} avec les candidats (emails personnalisés, planification de rendez-vous via calendrier).
\end{itemize}

L'ensemble de ces fonctionnalités repose sur des workflows intelligents, orchestrés par des outils d'automatisation, permettant de relier efficacement les différentes étapes du processus RH.


\section{Analyse des besoins}
\label{sec:analyse-besoins}

\subsection{Identification des acteurs}
\label{subsec:identification-acteurs}

Les principaux acteurs du système sont :

\begin{itemize}
    \item \textbf{Administrateur RH} : gère l'ensemble de la plateforme, les utilisateurs et les configurations;
    
    \item \textbf{Recruteur} : utilise la plateforme pour gérer les candidatures, évaluer les profils et communiquer avec les candidats;
    
    \item \textbf{Candidat} : accède au site web public pour consulter les offres et postuler;
    
    \item \textbf{Système intelligent (IA + workflows)} : automatise les tâches de traitement et d'analyse;
    
    \item \textbf{Systèmes externes} : calendrier, service d'email, authentification (Keycloak).
\end{itemize}


\subsection{Besoins fonctionnels}
\label{subsec:besoins-fonctionnels}

Le système doit permettre :

\begin{itemize}
    \item La gestion des utilisateurs et des rôles (admin, recruteur, candidat);
    
    \item L'importation et l'analyse automatique des CV;
    
    \item La modification et l'enrichissement des profils;
    
    \item L'interaction avec un chatbot intelligent (texte ou voix);
    
    \item La génération de CV standardisés selon différents templates;
    
    \item La gestion des offres d'emploi (création, modification, publication);
    
    \item La réception et le traitement des candidatures;
    
    \item Le classement automatique des candidats via scoring;
    
    \item L'envoi d'emails personnalisés;
    
    \item La planification automatique des entretiens;
    
    \item Le suivi et l'historique des candidatures.
\end{itemize}


\subsection{Besoins non-fonctionnels}
\label{subsec:besoins-non-fonctionnels}

Au-delà des fonctionnalités, le système doit respecter plusieurs contraintes techniques :

\begin{itemize}
    \item \textbf{Performance} : traitement rapide des CV et des workflows (temps de réponse < 2s);
    
    \item \textbf{Scalabilité} : possibilité d'ajouter des templates et gérer un grand volume de données;
    
    \item \textbf{Sécurité} : gestion des accès, authentification sécurisée, protection des données sensibles;
    
    \item \textbf{Fiabilité} : disponibilité continue du système (objectif de 99.5\%);
    
    \item \textbf{Maintenabilité} : architecture modulaire et code facilement maintenable;
    
    \item \textbf{Interopérabilité} : intégration avec différents services (API externes, calendrier, email).
\end{itemize}


\subsection{Diagramme de cas d'utilisation}
\label{subsec:diagramme-cas-utilisation}

La figure \ref{fig:use-case-diagram} présente les principaux cas d'utilisation du système. Ce diagramme montre les interactions entre les différents acteurs et les fonctionnalités proposées.

\begin{figure}[!ht]
    \centering
    \fbox{\parbox{\textwidth}{
        \textit{Remarque : Un diagramme UML détaillé des cas d'utilisation devrait être inséré ici, montrant les interactions entre les acteurs (Administrateur, Recruteur, Candidat) et les cas d'utilisation (Importer CV, Évaluer profil, Publier offre, etc.). Voir section \ref{subsec:diagramme-cas-utilisation}.}
    }}
    \caption{Diagramme de cas d'utilisation du système}
    \label{fig:use-case-diagram}
\end{figure}

Les principaux cas d'utilisation incluent :

\begin{enumerate}
    \item \textbf{Importer et analyser un CV} : le recruteur charge un document, le système l'analyse et en extrait les informations;
    
    \item \textbf{Transformer un CV} : le système applique un template pour standardiser la présentation;
    
    \item \textbf{Consulter et modifier un profil} : les utilisateurs peuvent visualiser et mettre à jour les données;
    
    \item \textbf{Publier une offre d'emploi} : création et publication d'une offre, synchronisation avec le site public;
    
    \item \textbf{Évaluer et scorer une candidature} : le système classe automatiquement les candidats;
    
    \item \textbf{Communiquer avec un candidat} : envoi d'emails automatisés, planification d'entretiens.
\end{enumerate}


\section{Méthodologie de développement}
\label{sec:methodologie}

\subsection{Choix de l'approche Agile}
\label{subsec:approche-agile}

Le projet intégrant des composants d'intelligence artificielle et des besoins évolutifs, une approche agile a été privilégiée.

Contrairement aux méthodes classiques (Waterfall), cette approche permet :

\begin{itemize}
    \item Une \textbf{adaptation continue} aux changements et aux découvertes;
    
    \item Une \textbf{validation progressive} des fonctionnalités avec les utilisateurs;
    
    \item Une \textbf{meilleure collaboration} entre l'équipe de développement et les responsables RH;
    
    \item Une \textbf{livraison incrémentale}, permettant d'utiliser des parties du système dès qu'elles sont prêtes.
\end{itemize}

L'IA, par nature expérimentale, bénéficie particulièrement de cette approche itérative.


\subsection{Méthode Scrum}
\label{subsec:methode-scrum}

Le développement est organisé en sprints de deux semaines. Les rôles principaux sont :

\begin{itemize}
    \item \textbf{Product Owner} : définit les priorités et valide les fonctionnalités;
    
    \item \textbf{Équipe de développement} : conçoit, développe et teste les fonctionnalités;
    
    \item \textbf{Scrum Master} : facilite les cérémonies Scrum et élimine les obstacles.
\end{itemize}

Chaque sprint comprend :

\begin{enumerate}
    \item Une \textbf{planification} (Sprint Planning) : définition des objectifs et des tâches;
    
    \item Des \textbf{mêlées quotidiennes} (Daily Standup) : synchronisation de l'équipe;
    
    \item Une \textbf{revue de sprint} (Sprint Review) : démonstration des fonctionnalités livrées;
    
    \item Une \textbf{rétrospective} (Sprint Retrospective) : analyse des améliorations possibles.
\end{enumerate}


\subsection{Organisation du projet}
\label{subsec:organisation-projet}

Le projet est structuré en plusieurs sprints couvrant :

\begin{itemize}
    \item \textbf{Sprint 0} : mise en place technique (backend, frontend, workflows, Docker);
    
    \item \textbf{Sprint 1-2} : gestion des profils et import de CV;
    
    \item \textbf{Sprint 3-4} : transformation et standardisation des CV;
    
    \item \textbf{Sprint 5-6} : gestion des candidatures et scoring;
    
    \item \textbf{Sprint 7-8} : intégration de l'intelligence artificielle et optimisations;
    
    \item \textbf{Sprint 9+} : automatisation des processus RH et déploiement.
\end{itemize}

Cette structure permet une progression logique du projet, avec des jalons clairement identifiés et des livrables testables à chaque étape.

\newpage
